%% Mouvement de rotation autour d'un axe fixe : hélice bipale tournant autour de O.
%% Chaque point de la pale décrit un cercle centré sur l'axe ; la norme du vecteur
%% vitesse croît proportionnellement au rayon (V = omega x R). omega et V dans le même sens (trigonométrique).
\documentclass[tikz,border=4pt]{standalone}
\usepackage{liaisons}
\begin{document}
\begin{tikzpicture}[x=1cm,y=1cm, scale=1.22]
\def\ang{10}   % inclinaison de la pale

% --- cercles décrits par les points B, C, D, E ------------------------
\foreach \r in {0.7,1.2,1.7,2.2}
  { \draw[solideA, dashed, dash pattern=on 2pt off 2pt, line width=0.5pt] (0,0) circle (\r); }

% --- l'hélice (deux pales) -------------------------------------------
\begin{scope}[rotate=\ang]
  \filldraw[fill=black!12, draw=black!60, line width=0.8pt]
     (2.45,0) to[out=100,in=8]   (0,0.30)
              to[out=188,in=80]  (-2.45,0)
              to[out=-80,in=-172](0,-0.30)
              to[out=-8,in=-100] cycle;
\end{scope}
\filldraw[fill=black!45, draw=black!70, line width=0.7pt] (0,0) circle (0.15);
\node[font=\small, below left=0pt] at (0,0) {$O$};

% --- points de la pale et vecteurs vitesse tangents -------------------
\foreach \r/\n in {0.7/B, 1.2/C, 1.7/D, 2.2/E} {
  \fill[solideA] (\ang:\r) circle (1.6pt);
  \draw[-{Stealth[length=5pt]}, line width=1.2pt, solideA]
       (\ang:\r) -- ++({\ang+90}:{0.34*\r});
  \node[font=\small, anchor=north, inner sep=2.5pt] at (\ang:\r) {$\n$};
}
\node[font=\small, text=solideA, anchor=south west] at (2.05,1.05) {$\vec{V}$};

% --- vitesse de rotation ----------------------------------------------
\draw[-{Stealth[length=5pt]}, line width=1pt, solideE]
     (-125:0.62) arc (-125:-55:0.62);
\node[font=\small, text=solideE] at (0,-0.86) {$\omega$};

% --- cote du rayon R (dans le repère de la pale) ----------------------
\begin{scope}[rotate=\ang]
  \draw[line width=0.4pt, black!70] (0,0.18) -- (0,0.95) (2.2,0.05) -- (2.2,0.95);
  \draw[{Stealth[length=4pt]}-{Stealth[length=4pt]}, line width=0.6pt, black!80] (0,0.85) -- (2.2,0.85);
  \node[font=\small, fill=white, inner sep=1pt] at (1.1,0.85) {$R$};
\end{scope}

% --- relation ---------------------------------------------------------
\node[font=\small, anchor=west] at (-2.6,-1.85) {$V = \omega \times R$};
\end{tikzpicture}
\end{document}
